%% ===== Project Summary (1 page) =====
%% NSF requires three labeled parts: Overview, Intellectual Merit, Broader Impacts.
%% The last line must carry 2-5 keywords (NSF 26-506, Section V.A).

\begin{center}
  {\large\bfseries Project Summary}\\[0.3em]
  {\large\bfseries PESOSE: Track 3: \sysname: AI-Enabled Digital Twins \\ 
  for Continuous Red Teaming of Open-Source Ecosystems}
\end{center}

\noindent\textbf{Overview}\\
Assessing the security posture of large computing infrastructures is
notoriously difficult. These are complex systems, dynamic in their
connectivity and usage patterns, and they often carry out processes that must
not be disrupted. Performing security analysis directly on a production
system can cause problems, as scans introduce overhead and analyzing services
for vulnerabilities can trigger crashes. Because of these risks, a security
assessment is usually performed manually, sporadically, and in an ad hoc
fashion by scarce human experts, while adversaries, now increasingly
augmented with AI, operate continuously. The goal of this proposal is the development of \sysname, an agentic framework
for the continuous, in-depth, and non-disruptive security assessment of
computing infrastructures built on open-source software. Such infrastructures
offer a unique opportunity for a new approach, because the openness of the
software enables the faithful replication of an infrastructure's topology and
service configuration. \sysname uses cooperating teams of AI agents to capture the configuration of
a target system, observe its execution, build a \emph{digital twin} of the
system, run red team exercises against the twin instead of the production
system, and, when vulnerabilities are found, guide a human-driven
verification and remediation of the issues in the original infrastructure.
Because the analysis runs on the twin, it can be repeated continuously and
adapted as the target evolves.

The ecosystem that this project targets is anchored on two existing, widely
deployed open-source products. The first is \textbf{nginx}, the reverse proxy
through which traffic enters and moves within nearly every open-source
deployment, and which we will harden against the memory-safety,
configuration, and request-smuggling defect classes that dominate its
history. The second is \textbf{ZMap}, the scanner with which operators
discover what is actually running, and which we will extend into an
instrument for safe, budgeted discovery inside production networks. Every
finding that we confirm on a twin is driven upstream to the project
responsible, and the work therefore improves not only individual deployments,
but also the shared ecosystem on which they are built.

\vspace{0.2ex}
\noindent\textbf{Intellectual Merit}\\
The \sysname framework contributes new agentic workflows that automate the
tedious tasks of capturing an infrastructure's configuration and reproducing
it in a virtualized environment. These workflows are supported by novel
techniques for replicating large-scale infrastructure at reduced scale while
preserving its security-relevant properties, and for measuring the
representativeness of the resulting twin through logic-based invariants and
attack-graph preservation analysis. A distinguishing feature of the approach
is that the fidelity of the twin is quantified rather than assumed, and that
the soundness of an analysis is reported together with a measure of how much
of the target system the capture process actually resolved. Vulnerability-analysis agents then examine the twin using novel techniques
for chaining vulnerabilities across a networked infrastructure, extending the
investigators' prior work on AI-planning-based privilege-escalation chains
and on a universal theory of in-program exploitation to complete
open-source ecosystems.

\vspace{0.2ex}
\noindent\textbf{Broader Impacts}\\
Open-source software underpins the Nation's critical infrastructure, from
health care and financial systems to utilities and government services. By
enabling continuous, non-disruptive security assessment of the ecosystems
built on this software, \sysname will improve the resiliency of that
infrastructure and help prevent attacks with catastrophic consequences. The improvements that the project contributes to nginx and ZMap
reach every downstream user of those products, including the many
organizations that cannot afford commercial alternatives. The project will
also produce open-source artifacts and a public corpus of reference
infrastructures, integrate infrastructure-security modules into the curricula
at UC Santa Barbara and Arizona State University, involve undergraduate students in research, and organize a new public capture-the-flag
competition centered on the creation of digital twins.

\vspace{0.2ex}
\noindent\textbf{Keywords:} critical infrastructure; open-source security; digital twins; AI agents; vulnerability analysis.
